\subsection{Main Result: Upper Bound on Excess Risk}

We now state our main theorem, which provides an explicit upper bound on the excess risk as a function of the sample sizes $m$ and $n$, relevant distribution-dependent quantities, and the choice of regularizer $\lamz$.

Let
\begin{equation}\label{eqn: uniform weight bound}
    \mathsf{W} := \max\{\sup_\lambda\|\wspsl\|,\sup_\lambda\|\whl\|,\sup_\lambda\|\wsl\|,\|\ws\|,\|\us\|,\|\uh\|\}.
\end{equation}
We show in Proposition \ref{cor: uniform bound on weights} that with high probability, $\mathsf{W} < \infty$.


\begin{thm}\label{thm: main theorem}
    Let $\us$ and $\vs$ be as defined in \eqref{eqn: u star} and \eqref{eqn: v star}, respectively. Let $\ws$, $\wsn$, and $\mathsf{W}$ be as defined in \eqref{eqn: w star}, \eqref{eqn: w lambda star} and \eqref{eqn: uniform weight bound} respetively.
Consider the partitioned training datasets $D = \{(x_i, y_i)\}_{i=1}^n$ and $S = \{(x_j, z_j)\}_{j=1}^{m}$, as defined in \eqref{eqn: partioned dataset}, and let $\delta > 0$. Take the regularizer $\lambda$ as
\begin{equation*}
    \lambda = \left(\frac{1}{\sqrt{n}}, \frac{1}{\sqrt{n}}, \lamz\right) \in \mathbb{R}^{3},
\end{equation*}
where $\lamz \geq \frac{1}{\sqrt{n}}$, controls the regularization strength for the side information component.
Then, with probability at least $1-4\delta$, the following inequality holds.
   \begin{equation}\label{eqn: main equation}
        \begin{aligned}
            R(\whl) - R(\ws) &\leq \frac{C_1}{\sqrt{n}} +
            \min\Bigg\{\frac{\lamz|\wsz|^2}{T(\vs) - R(\ws)},
            1+\frac{\lamz |\wslz|^2}{T(\vs) - R(\ws)}\Bigg\}\Big(T(\vs) - R(\ws)\Big) \\
            &+
            \min\Bigg\{\frac{\ln(2)}{\lamz},\mathsf{W}\Bigg\}\ln(4)\Bigg(\frac{C_2}{\sqrt{m}} + L(\us)\Bigg),
        \end{aligned}  
    \end{equation}
  where $C_1 = \mathcal{O}\left( W^2 + \kappa W \sqrt{\ln(1/\delta)} + \kappa^2 \ln(1/\delta) \right)$ and $C_2 = \mathcal{O}\left( \|u^*\| \sqrt{\ln(1/\delta)} \right)$ are constants.\footnote{The consequence of the assmption $m < n$ is that the term $1/\sqrt{m}$ on the RHS does not get absorbed in the $1/\sqrt{n}$.}
\end{thm}


\begin{proof}[Proof Outline]
To analyze the excess risk, we decompose the left-hand side of \eqref{eqn: main equation} into four terms, each capturing a distinct source of error:
\begin{equation*}
\begin{aligned}
    R(\whl) - R(\ws) =\;&
    \underbrace{R(\whl) - \rph(\whl)}_{\mathrm{(A)}} +
    \underbrace{\rph(\whl) - \rph(\wsl)}_{\mathrm{(B)}} \\
    &+ \underbrace{\rph(\wsl) - R(\wsl)}_{\mathrm{(C)}} +
    \underbrace{R(\wsl) - R(\ws)}_{\mathrm{(D)}}.
\end{aligned}
\end{equation*}

\noindent The first and third terms, (A) and (C), quantify the estimation error resulting from the use of the imputed side information $\ph$. This error depends on the complexity of the relationship between $(X, Z)$ quantified by $L(\us)$, the sample size $m$ of the side information dataset $S$, and the regularization parameter $\lamz$, which determines the extent to which the model relies on the imputed side information. These terms contribute the following to the
upper bound:
\begin{equation*}
    \min\Bigg( \frac{\ln(2)}{\lamz},\ \mathsf{W} \Bigg)
    \times
    \ln(4)\Bigg( \frac{C_2}{\sqrt{m}} + L(\us) \Bigg).
\end{equation*}

\noindent The second term, (B), corresponds to the estimation error in the standard supervised learning setting for the primary classifier, which results in the contribution $    \frac{C_1}{\sqrt{n}}$
to the upper bound in \eqref{eqn: main equation}.

\noindent Finally, the fourth term, (D), captures the gap between the risk
of the optimal regularized classifier and the risk of the optimal classifier. This term contributes
\begin{equation*}
\begin{aligned}
    &\min\left\{\frac{\lamz |\ws_z|^2}{T(\vs) - R(\ws)},\
    1 + \frac{\lamz |\ws_{\lamz}|^2}{T(\vs) - R(\ws)}\right\}\left(T(\vs) - R(\ws)\right)
\end{aligned}
\end{equation*}
to the upper bound in \eqref{eqn: main equation}.
Note that $\vs$ coincides with $\wsl$ when $\lamz = \infty$, so this term quantifies the proportion of the worst-case gap $T(\vs) - R(\ws)$ (with the worst at $\lamz = \infty$) that is realized, depending on the choice of the regularizer $\lamz$.
\end{proof}

\subsubsection{Interpretation of the Excess Risk Bound}

To provide intuition for the upper bound in Theorem \ref{thm: main theorem}, consider the university admission scenario. Here, the feature $x \in \R^d$ represents a candidate's academic scores, the side information $z$ is the admission decision of a secondary university, and the label $y$ is the decision of the primary university. The excess risk $R(\hat{w}_{\lambda}) - R(w^*)$ represents the performance gap between our learned classifier $\hat{w}_{\lambda}$ and the benchmark optimal classifier $w^*$, which has access to the true joint distribution of $(X, Y, Z)$. This risk is controlled by the regularization parameter $\lambda_z$ and consists of three distinct components.

The first term, $\mathcal{O}(1/\sqrt{n})$, is the standard estimation error associated with the sample size $n$ of the primary university's dataset. This error is unavoidable and decays as more primary data becomes available.

The second term captures the potential gain from incorporating $z$. It is proportional to the gap $T(v^*) - R(w^*)$, which measures how much knowing the true side information improves prediction accuracy over using features $x$ alone. This term is controlled linearly by $\lambda_z$. If $z$ carries meaningful information (i.e., the gap is significant), the bound is lowered by reducing $\lambda_z$, which encourages the model to assign larger weight to the side information.

The third term quantifies the error introduced by imputing $z$ from $x$ using the secondary university sample. It depends on the sample size $m$ of the secondary university's data and the optimal prediction error $L(u^*)$, which reflects the complexity of the relationship between features and side information. This term is inversely controlled by $\lambda_z$ (scaling as $1/\lambda_z$). If the features cannot accurately predict $z$, the imputed values are noisy, and a large $\lambda_z$ is required to ignore the unreliable imputed data.

The interplay between the second and third terms creates a tension in selecting $\lambda_z$. If $z$ is highly informative for predicting $y$ (large gap $T(v^*) - R(w^*)$), a smaller $\lambda_z$ is desirable to leverage this information. However, if $x$ poorly predicts $z$ (large $L(u^*)$), a larger $\lambda_z$ is necessary to mitigate the noise from imputed side information. The regularizer $\lambda_z$ must balance these two competing terms. In this case, the optimal $\lambda_z$ must find a balance to extract the signal in $z$ while damping the noise from imperfect imputation. In most other cases, where $z$ is either useless or perfectly predictable, the terms do not conflict, and the choice of $\lambda_z$ becomes straightforward. As we will show in subsequent sections, despite this tension, the optimal $\lambda_z$ often exhibits bang-bang behavior, taking on extreme values (either very large or very small) depending on which factor is more limiting: the noise introduced by imperfectly imputed $z$, or the dependence of $y$ on $z$. This distinction is determined by the underlying distributional properties.

\subsubsection{Retrospective Regularizer Selection}

Theorem \ref{thm: main theorem} provides an explicit upper bound on the excess risk as a function of the regularization parameter $\lamz$. This upper bound holds for all values of $\lamz$ and all distributions. We are interested in determining the value of $\lamz$ that minimizes this upper bound. Since $\lamz$ is a hyperparameter, its optimal value depends on the underlying population distribution. Under certain assumptions, we show that the choice for $\lamz$ may be restricted to either $0$ or $\infty$.
To facilitate this analysis, we work with a relaxed version of the upper bound \eqref{eqn: main equation} from Theorem \ref{thm: main theorem}. Specifically, by applying Lemma \ref{lem: bound weights in terms of lambda} (given in the Supplementary Proofs and Results), we have $\lamz |\wslz|^2 \leq \ln(2)$, so we obtain the following relaxed bound.




\begin{equation}\label{eqn: weaker main equation}
    \begin{aligned}
        R(\whl) - R(\ws) &\leq \frac{C_1}{\sqrt{n}} +
        \min\Big\{\frac{\lamz|\wsz|^2}{T(\vs) - R(\ws)},1+\frac{\ln(2)}{T(\vs) - R(\ws)}\Big\}\Big(T(\vs) - R(\ws)\Big) \\
        &+
        \min\Big\{\frac{\ln(2)}{\lamz},\mathsf{W}\Big\}\ln(4)\Big(\frac{C_2}{\sqrt{m}} + L(\us)\Big),
    \end{aligned}
\end{equation}
where $C_1 = \mathcal{O}\left( W^2 + \kappa W \sqrt{\ln(1/\delta)} + \kappa^2 \ln(1/\delta) \right)$ and $C_2 = \mathcal{O}\left( \|u^*\| \sqrt{\ln(1/\delta)} \right)$ are constants. This form makes the $\lamz$ dependence on excess risk explicit, enabling a more tractable analysis of its optimal value.

Let
\begin{equation*}\label{eqn: Optimal Lambda}
    \begin{aligned}
            \Lambda_z^\star :=& \arg\min_{\lamz}        \frac{C_1}{\sqrt{n}} +
        \min\Big\{\frac{\lamz|\wsz|^2}{T(\vs) - R(\ws)}, 1+\frac{\ln(2)}{T(\vs) - R(\ws)}\Big\}\Big(T(\vs) - R(\ws)\Big) \\
        &+
        \min\Big\{\frac{\ln(2)}{\lamz},\mathsf{W}\Big\}\ln(4)\Big(\frac{C_2}{\sqrt{m}} + L(\us)\Big),
    \end{aligned}
\end{equation*}
be the optimal (population distribution dependent) regularizer that minimizes the upper bound on the excess risk.
For a given value of $\Lambda_z^\star$, three disjoint and exhaustive cases arise, each corresponding to a different regime for the optimal regularizer:
\begin{itemize}
    \item \textbf{Case (a):}
    \begin{equation*}
        \frac{\Lambda_z^\star|\wsz|^2}{T(\vs) - R(\ws)} \geq 1+\frac{\ln(2)}{T(\vs) - R(\ws)}.
    \end{equation*}
    This is the regime considered in Corollary \ref{cor: consequence of main theorem 1}.
    \item \textbf{Case (b):}     \begin{equation*}
        \frac{\Lambda_z^\star|\wsz|^2}{T(\vs) - R(\ws)} < 1+\frac{\ln(2)}{T(\vs) - R(\ws)}.
    \end{equation*} and
    \begin{equation*}
        \frac{\ln(2)}{\Lambda_z^\star} \geq \mathsf{W}.
    \end{equation*} This is the regime considered in Corollary \ref{cor: consequence of main theorem 2}.
    \item \textbf{Case (c):}
    \begin{equation*}
        \frac{\Lambda_z^\star|\wsz|^2}{T(\vs) - R(\ws)} < 1+\frac{\ln(2)}{T(\vs) - R(\ws)}.
    \end{equation*} and
    \begin{equation*}
        \frac{\ln(2)}{\Lambda_z^\star} < \mathsf{W}.
    \end{equation*}
    This is the regime considered in Corollary \ref{cor: consequence of main theorem 3}.
\end{itemize}

Each case is formalized in the following corollaries, which provide guidance for selecting the regularization parameter in different regimes.

The first corollary shows that if $\Lambda_z^\star$ is sufficiently large, then it is $\Lambda_z^\star = \infty$.

\begin{corollary}[Regime I: Large Regularizer]\label{cor: consequence of main theorem 1}
Suppose
\begin{equation*}
    \frac{\Lambda_z^\star|\ws_z|^2}{T(\vs) - R(\ws)} \geq 1+\frac{\ln(2)}{T(\vs) - R(\ws)}.
\end{equation*}
Then $\Lambda_z^\star = \infty$ and
    \begin{equation*}
        R(\whl) - R(\ws) = \frac{C_1}{\sqrt{n}} + T(\vs) - R(\ws).
    \end{equation*}
\end{corollary}

\begin{proof}
Consider the condition
\begin{equation*}
    \frac{\Lambda_z^\star |\ws_z|^2}{T(\vs) - R(\ws)} \geq 1 + \frac{\ln(2)}{T(\vs) - R(\ws)}.
\end{equation*}
Under this condition, the minimum in the second term of the upper bound in \eqref{eqn: weaker main equation} is achieved by the second argument.

Now consider $\lamz$ such that $\lamz > \Lambda_z^\star$. This choice does not increase the contribution of the second term, as
\begin{equation*}
    \frac{\lamz |\ws_z|^2}{T(\vs) - R(\ws)} \geq \frac{\Lambda_z^\star |\ws_z|^2}{T(\vs) - R(\ws)} \geq 1 + \frac{\ln(2)}{T(\vs) - R(\ws)},
\end{equation*}
but it decreases the contribution of the third term in \eqref{eqn: weaker main equation} to $0$, as the first argument of the third term goes to zero. This decreases the value of the overall upper bound.

However, since $\Lambda_z^\star$ is assumed to be optimal, increasing $\lamz$ beyond it should not reduce the upper bound. Thus, $\Lambda_z^\star = \infty$. Substituting this value back into \eqref{eqn: main equation} and noting that $\lamz |\wslz|^2 \to 0$ as $\lamz \to \infty$, we obtain
    \begin{equation*}
        R(\whl) - R(\ws) = \frac{C_1}{\sqrt{n}} + T(\vs) - R(\ws).
    \end{equation*}
\end{proof}

The second corollary establishes that if $\Lambda_z^\star$ is sufficiently small, then the optimal value is $\Lambda_z^\star = \frac{1}{\sqrt{n}}$.

\begin{corollary}[Regime II: Small Regularizer]\label{cor: consequence of main theorem 2}
Suppose
    \begin{equation*}
        \frac{\Lambda_z^\star|\wsz|^2}{T(\vs) - R(\ws)} < 1+\frac{\ln(2)}{T(\vs) - R(\ws)}.
    \end{equation*}
and
    \begin{equation*}
        \frac{\ln(2)}{\Lambda_z^\star} \geq \mathsf{W}.
    \end{equation*}
Then $\Lambda_z^\star = \frac{1}{\sqrt{n}}$ and
    \begin{equation*}
    \begin{aligned}
        R(\whl) - R(\ws) &\leq \frac{C_1}{\sqrt{n}} + \mathsf{W}\ln(4)\Big(\frac{C_2}{\sqrt{m}} + L(\us)\Big).
        \end{aligned}
    \end{equation*}
\end{corollary}

\begin{proof}
Consider the condition
    \begin{equation*}
        \frac{\Lambda_z^\star|\wsz|^2}{T(\vs) - R(\ws)} < 1+\frac{\ln(2)}{T(\vs) - R(\ws)}.
    \end{equation*}
and
    \begin{equation*}
        \frac{\ln(2)}{\Lambda_z^\star} \geq \mathsf{W}.
    \end{equation*}
Under these conditions, the minimum in the second term and the third term of the upper bound in \eqref{eqn: weaker main equation} is achieved by the first argument, and the second argument respectively.

Let $\lamz$ be such that $\lamz < \Lambda_z^\star$.
This choice does not increase the contribution of the third term, as \begin{equation*}
        \frac{\ln(2)}{\lamz} \geq \frac{\ln(2)}{\Lambda_z^\star} \geq \mathsf{W},
    \end{equation*}
but it decreases the contribution of the second term as
    \begin{equation*}
        \frac{\lamz|\wsz|^2}{T(\vs) - R(\ws)}< \frac{\Lambda_z^\star|\wsz|^2}{T(\vs) - R(\ws)} < 1+\frac{\ln(2)}{T(\vs) - R(\ws)}.
    \end{equation*}
This decreases the value of the overall upper bound.

However, since $\Lambda_z^\star$ is assumed to be optimal, decreasing $\lamz$ below it should not reduce the upper bound. Thus, $\Lambda_z^\star = \frac{1}{\sqrt{n}}$, the minimum allowed value of $\lamz$. Substituting this value back into \eqref{eqn: main equation}, we obtain
    \begin{equation*}
    \begin{aligned}
        R(\whl) - R(\ws) &\leq \frac{C_1}{\sqrt{n}} + \mathsf{W}\ln(4)\Big(\frac{C_2}{\sqrt{m}} + L(\us)\Big).
        \end{aligned}
    \end{equation*}

\end{proof}

The third corollary shows that if $\Lambda_z^\star$ is neither a big value nor a small value, then it scales with the size of side information dataset $S$  as $\Lambda_z^\star = \mathcal{O}(m^{-1/4})$.

\begin{corollary}[Regime III: Intermediate Regularizer]\label{cor: consequence of main theorem 3}
Suppose
    \begin{equation*}
        \begin{aligned}
            \frac{\Lambda_z^\star|\wsz|^2}{T(\vs) - R(\ws)} &< 1+\frac{\ln(2)}{T(\vs) - R(\ws)} \quad \text{and} \\
            \frac{\ln(2)}{\Lambda_z^\star} &< \mathsf{W}.
        \end{aligned}
    \end{equation*}
Then the optimal regularizer satisfies
 \begin{equation*}
        \Lambda_z^\star = \frac{\sqrt{2(\ln(2))^2\Big(\frac{C_2}{\sqrt{m}} + L(\us)\Big)}}{|\wsz|},
    \end{equation*}
and the excess risk is bounded by
    \begin{equation*}
    \begin{aligned}
        R(\whl) - R(\ws) &\leq \frac{C_1}{\sqrt{n}} + 2|\wsz|\sqrt{2(\ln(2))^2\Big(\frac{C_2}{\sqrt{m}} + L(\us)\Big)}.
    \end{aligned}
    \end{equation*}
\end{corollary}

\begin{proof}
Suppose the following two conditions hold:
\begin{equation*}
    \begin{aligned}
        \frac{\Lambda_z^\star |\wsz|^2}{T(\vs) - R(\ws)} &\leq 1 + \frac{\Lambda_z^\star |\wslz|^2}{T(\vs) - R(\ws)}, \\
        \frac{\ln(2)}{\Lambda_z^\star} &\leq \mathsf{W}.
    \end{aligned}
\end{equation*}
Under these conditions, the minimum in both the second and third terms of the upper bound in \eqref{eqn: weaker main equation} is achieved by the first argument. Since the upper bound is minimized at $\Lambda_z^\star$, the derivative with respect to $\lambda_z$ must vanish at this point. This gives
\begin{equation*}
    |\wsz|^2 - \frac{\ln(2)}{\lamz^2} \ln(4)\Big(\frac{C_2}{\sqrt{m}} + L(\us)\Big)\Bigg|_{\lamz = \Lambda_z^\star} = 0.
\end{equation*}
Solving this yields
\begin{equation*}
    \Lambda_z^\star = \frac{\sqrt{2(\ln(2))^2\Big(\frac{C_2}{\sqrt{m}} + L(\us)\Big)}}{|\wsz|},
\end{equation*}
Substituting this value back into \eqref{eqn: main equation} gives
\begin{equation*}
    R(\whl) - R(\ws) \leq \frac{C_1}{\sqrt{n}} + 2|\wsz| \sqrt{ 2(\ln(2))^2 \left( \frac{C_2}{\sqrt{m}} + L(\us) \right) }.
\end{equation*}
\end{proof}

We now show that the Regime III considered in Corollary \ref{cor: consequence of main theorem 3} does not arise.

\begin{corollary}
Let $n$ and $m$ be sufficiently large. Then the optimal value of $\Lambda^\star_z$ is either $\infty$ or $0$.
\end{corollary}

\begin{proof}
If
\begin{equation*}
    \frac{\Lambda_z^\star|\ws_z|^2}{T(\vs) - R(\ws)} \geq 1+\frac{\ln(2)}{T(\vs) - R(\ws)},
\end{equation*}
then from Corollary \ref{cor: consequence of main theorem 1}, we have $\Lambda^\star_z = \infty$.
If
\begin{equation*}
    \frac{\Lambda_z^\star|\ws_z|^2}{T(\vs) - R(\ws)} < 1+\frac{\ln(2)}{T(\vs) - R(\ws)},
\end{equation*}
then, according to Corollaries \ref{cor: consequence of main theorem 2} and \ref{cor: consequence of main theorem 3}, there are two possible values for the optimal regularizer.
If $\Lambda^\star_z = \frac{1}{\sqrt{n}} \to 0$ (Corollary \ref{cor: consequence of main theorem 2}), the upper bound on the excess risk is
\begin{equation}\label{eqn: upperbound cor 2}
\begin{aligned}
    R(\whl) - R(\ws) &\leq \frac{C_1}{\sqrt{n}} + \mathsf{W}\ln(4)\Big(\frac{C_2}{\sqrt{m}} + L(\us)\Big)
\end{aligned}
\end{equation}
and if $\Lambda_z^\star = \frac{\sqrt{2(\ln(2))^2\left(\frac{C_2}{\sqrt{m}} +L(\us)\right)}}{|\ws_z|}$(Corollary \ref{cor: consequence of main theorem 3}), the upper bound on the excess risk is
\begin{equation}\label{eqn: upperbound cor 3}
    R(\whl) - R(\ws) \leq \frac{C_1}{\sqrt{n}} + 2|\ws_z|\sqrt{2(\ln(2))^2\Big(\frac{C_2}{\sqrt{m}} +L(\us)\Big)}.
\end{equation}
Note that the logistic risk of the zero predictor is $L(0) = \ln(2) < 1$. By the optimality of $\us$, we have $L(\us) \leq L(0) < 1$. The bound in \eqref{eqn: upperbound cor 2} scales as $\mathcal{O}\left(\frac{1}{\sqrt{m}} + L(\us)\right)$, whereas the bound in \eqref{eqn: upperbound cor 3} scales as $\mathcal{O}\left(\frac{1}{m^{1/4}} + \sqrt{L(\us)}\right)$. Thus, for sufficiently large $n$ and $m$, since $L(\us) < 1$, the bound in \eqref{eqn: upperbound cor 2} is smaller, implying Regime III considered in Corollary \ref{cor: consequence of main theorem 3} will never occur, and hence $\Lambda^\star_z = 0$.



\end{proof}


% \subsubsection{Examples Where Upper Bound Is Asymptotically Tight}

% We now provide an example of a distribution for which the upper bound in Theorem \ref{thm: main theorem} is asymptotically tight.

% \begin{eg}
% Let $(X,Y)$ be any distribution on $\cg{X} \times \cg{Y}$. Now consider the distribution $(X,Y,Z)$ defined by:
% \begin{equation*}
%     \begin{array}{|c|c|}
%         \hline
%         \text{Event} & \text{Probability} \\
%         \hline
%         \prob(Z=+1|X=x,Y=y) & 1\\
%         \hline
%         \prob(Z=-1|X=x,Y=y) & 0 \\
%         \hline
%     \end{array}
% \end{equation*}
% From the structure of this distribution, it follows that $\ws = \vs$ (and hence $R(\ws) = T(\vs)$). Furthermore, by setting $\lamz = \infty$ and applying Theorem \ref{thm: main theorem}, we obtain:
% \begin{equation*}
%     \begin{aligned}
%         R(\whl) - R(\ws) &= L(\vh) - T(\vs)\\
%         &\leq \cg{O}\left(\frac{1}{\sqrt{n}}\right) + T(\vs) - R(\ws)\\
%         &= \cg{O}\left(\frac{1}{\sqrt{n}}\right).
%     \end{aligned}
% \end{equation*}

% This demonstrates that the upper bound in Theorem \ref{thm: main theorem} is asymptotically tight for this distribution.
% \end{eg}